\subsection{Методы снижения размерности} \label{sec:dimred_methods}

Задача снижения размерности состоит в нахождении представления многомерных данных в пространстве существенно меньшей размерности при сохранении наиболее значимых закономерностей исходного набора.
В нейронауке эта задача приобретает особое значение в связи с быстрым ростом числа одновременно регистрируемых нейронов \cite{Stevenson2011}: когда размерность пространства состояний популяции составляет сотни и тысячи измерений, прямая интерпретация данных становится невозможной, и методы снижения размерности выступают основным инструментом извлечения коллективных переменных \cite{Cunningham2014}.

\paragraph{Линейные методы.}
Наиболее распространённым линейным подходом является \textit{анализ главных компонент} (PCA), проецирующий данные на ортогональные направления максимальной дисперсии.
PCA широко используется для визуализации нейронных траекторий \cite{Gallego2017}, а также для предварительной оценки размерности пространства активности~--- например, числа главных компонент, необходимых для объяснения заданной доли дисперсии \cite{Cunningham2014}.
К линейным методам относятся также \textit{линейный дискриминантный анализ} (LDA), ищущий проекции с максимальным разделением заранее заданных классов, и \textit{многомерное масштабирование} (MDS), сохраняющее попарные расстояния между точками.
Общее ограничение линейных методов заключается в том, что они связывают низкоразмерное представление с исходными данными линейным преобразованием.
Если данные лежат вблизи линейного подпространства, такие методы восстанавливают их структуру адекватно; однако нейронные многообразия зачастую существенно нелинейны (секция~\ref{sec:population_coding}), и линейные проекции могут сильно искажать геометрию активности.

\paragraph{Нелинейные методы: обучение на многообразиях.}
Для восстановления нелинейной структуры данных разработано семейство методов, объединяемых под общим названием \textit{обучение на многообразиях} (manifold learning).
Эти алгоритмы исходят из предположения, что данные, вложенные в высокоразмерное пространство, на самом деле лежат вблизи гладкого многообразия существенно меньшей размерности, и используют локальную геометрию данных для построения низкоразмерного представления.

\textit{Isomap} \cite{tenenbaum2000}~--- один из первых методов этого семейства. Алгоритм строит граф ближайших соседей в исходном пространстве, аппроксимирует геодезические расстояния между всеми парами точек через кратчайшие пути в графе, а затем применяет классическое MDS к полученной матрице расстояний.
В контексте нейроданных Isomap использовался, в частности, для построения низкоразмерных вкраплений функциональных связей по данным фМРТ состояния покоя \cite{gallos2021}.

\textit{Локально линейное вложение} (LLE) \cite{Roweis2000} восстанавливает глобальную структуру многообразия на основе локальных линейных приближений: каждая точка выражается как линейная комбинация своих ближайших соседей, после чего ищется низкоразмерная конфигурация, в которой эти локальные линейные зависимости сохраняются.
LLE применялся к временным рядам фМРТ для нелинейного извлечения информативных признаков \cite{mannfolk2010, Li2016}.

\textit{Лапласовы собственные карты} (Laplacian Eigenmaps, LE) \cite{Belkin2003} основаны на спектральном разложении графового лапласиана, построенного по матрице близости между точками.
Метод оптимизирует вложение таким образом, чтобы близкие в исходном пространстве точки оставались близкими и в проекции.
LE использовались для кластеризации нейронной активности, классификации состояний мозга по данным фМРТ \cite{Pospelov2021, Khosla2019} и для анализа сигналов ЭЭГ \cite{Ataee2007}.

\textit{Диффузионные карты} (diffusion maps) \cite{coifman2006} обобщают спектральный подход, используя вероятности перехода случайного блуждания на графе данных.
Евклидово расстояние в результирующем пространстве приближает так называемое диффузионное расстояние, которое учитывает связность данных в целом, а не только локальные евклидовы расстояния.
В нейровизуализации диффузионные карты применялись для разделения активированного и фонового многообразий в событийно-связанных данных фМРТ \cite{shen2005}.

Общими ограничениями классических методов обучения на многообразиях являются: (1)~вычислительная сложность порядка $O(N^a)$, $2 < a < 3$, связанная со спектральным разложением матриц аффинности, что затрудняет применение к большим наборам данных; (2)~отсутствие естественного способа вычислить вложение для новых точек, не входивших в обучающую выборку (так называемая проблема out-of-sample).

\paragraph{Современные методы: t-SNE и UMAP.}
Альтернативный подход к нелинейному снижению размерности основан на оптимизации некоторого функционала расхождения между распределениями расстояний в исходном и целевом пространствах.

\textit{t-SNE} \cite{vandermaaten08a} строит распределения вероятностей на парах точек в высокоразмерном пространстве (гауссово ядро) и в низкоразмерном пространстве (ядро Стьюдента), после чего минимизирует расхождение Кульбака--Лейблера между этими распределениями.
Метод хорошо выявляет кластерную структуру данных и активно применяется для визуализации нейронных популяций, однако его результаты чувствительны к выбору гиперпараметров, а глобальные расстояния между кластерами могут сильно искажаться.

\textit{UMAP} \cite{McInnes2018} основан на теории нечётких симплициальных множеств и также оптимизирует сходство локальных структур в двух пространствах, однако, в отличие от t-SNE, лучше сохраняет глобальную структуру данных, масштабируется на существенно большие выборки и допускает вложение новых точек.
В нейронауке UMAP получил широкое распространение как инструмент визуализации и кластеризации многомерной нейронной активности.

Следует, однако, отметить, что и t-SNE, и UMAP ориентированы прежде всего на визуализацию в двух--трёх измерениях и не гарантируют сохранение метрических свойств данных, что ограничивает их использование для количественного анализа геометрии нейронных многообразий.

\paragraph{Автокодировщики.}
Принципиально иной подход к нелинейному снижению размерности реализуют \textit{автокодировщики} (autoencoders)~--- нейросетевые архитектуры, обучающиеся сжимать входные данные в низкоразмерное латентное представление и затем восстанавливать исходный сигнал из него.
В отличие от методов обучения на многообразиях, автокодировщики обучают параметрическое отображение, которое может быть непосредственно применено к новым данным, и масштабируются на существенно большие выборки.

Одним из наиболее влиятельных примеров применения автокодировщиков к нейронным данным является модель LFADS~\cite{Pandarinath2018}, подробно рассмотренная в следующей секции.

Таким образом, автокодировщики занимают промежуточное положение между классическими методами обучения на многообразиях и полноценными моделями латентной динамики, обсуждаемыми далее, и представляют собой гибкий инструмент нелинейного снижения размерности нейронных данных.
